\section{Related Works}

\paragraph{SST with LLM} Recent work shows that simultaneous speech translation (SST) can achieve markedly higher translation quality when large language models (LLMs) serve as the backbone~\citep{ahmad-etal-2024-findings}. \citet{koshkin-etal-2024-transllama} adapt an LLM to SST by finetuning on a small set of synthetic translation trajectories. \citet{ouyang-etal-2025-infinisst} further scale this idea with an interleaving architecture that supports efficient inference over unbounded speech streams. \citet{fu2025efficientadaptivesimultaneousspeech} study similar architectures and propose alternative training strategies and data-synthesis pipelines. \citet{guo2025streamuniachievingstreamingspeech} propose a unified model for streaming transcription and translation, introducing a truncation mechanism that prunes historical speech and text based on automatically transcribed inputs. 
However, most existing SST systems still struggle with domain-specific terminology even using LLMs, whereas \method integrates a terminology retriever into the SST process to translate specialized terms more accurately and reliably.
% \textcolor{red}{check latest works}

\paragraph{Retrieval-Augmented Translation}

Translating rare words and domain-specific terminology remains a persistent challenge for modern translation systems. A broad line of work addresses this via non-parametric retrieval at inference time: kNN-MT augments neural MT with nearest-neighbor lookup over a large datastore to improve domain adaptation and rare-token translation~\citep{khandelwal2021nearest}, and follow-up methods improve robustness by learning kernel-smoothed MT with retrieval~\citep{jiang-etal-2021-learning} as well as improving efficiency (e.g., faster and memory-efficient retrieval pipelines)~\citep{meng-etal-2022-fast,zhu-etal-2023-ink}. In parallel, recent LLM-based MT studies investigate incorporating domain knowledge at inference time through retrieval (demonstrations or terminologies)~\citep{li-etal-2025-leveraging}. Closer to speech translation, \citet{li-etal-2024-optimizing-rare} propose a retrieval-and-demonstration framework that prepends retrieved speech segments as demonstrations for offline ST, and \citet{wu-etal-2025-locate} introduce sliding-window retrieval to localize terminology occurrences and inject retrieved terms into a speech-LLM prompt. However, these approaches are primarily designed for offline translation, whereas \method targets streaming translation and integrates retrieval directly into the incremental translation process.

\paragraph{Contextual-Biasing ASR}
The recognition of rare words and domain-specific terminology, often referred to as contextual biasing, has also been extensively studied in automatic speech recognition (ASR). Recent work explores injecting bias lists directly into Speech LLM prompts ~\citep{10446898,gong25_interspeech}. Contextual-biasing has also been investigated in streaming settings~\citep{xu23d_interspeech}, but not with speech LLMs. 